% Transcription — fonds Alexandre Grothendieck, shelfmark 1, batch 1 (pages 1-20).
% Established from the facsimile archives/batches/1/batch-01.pdf (a mirror of
% https://grothendieck.umontpellier.fr/1.pdf), per the skill
% transcribe-grothendieck.
%
% Pass: Fable 5.1 (claude-fable-5-1), 2026-09-05 — first pass, unchecked against
% the pages by a human.
%
% Nineteen of the twenty pages are transcribed. Page 1 is blank and absent.
%
% \page{} numbers the position in the folder — what the inventory counts
% (149 pages) and what the site's facsimile pane indexes. The sheets carry the
% archivists' pencil numbers, and from the third sheet on those run two ahead:
% sheet 3 is pencilled 5, sheet 20 is pencilled 22 (pencil 3 and 4 have no
% scan). A note at page 3 records this, as folder 134-1 did for the same
% situation.
%
% The batch falls into three runs, and two hands' worth of legibility.
%
%   Page 2 — his plan of the notes in seven numbered sections, twice
%     renumbered: réarrangements de fonctions, fonctions de Weyl, inégalités,
%     mesures spectrales d'un opérateur, compléments sur le déterminant, les
%     inégalités pour AB, autres inégalités.
%   Pages 3-5 — a separate fragment on folded sheets, in a fast blue hand: the
%     left multipliers Γ_g of a C*-algebra C inside its bidual C'', the
%     two-sided multipliers Γ, and the isomorphism φ : Γ_1 → Γ with inverse ψ
%     when C is embedded in a von Neumann algebra R, proved through the
%     filter of approximate units. Ends "CQFD". The formulas carry; most of
%     the connecting prose is \ill{} or \uncertain{}.
%   Pages 6-20 — the notes the plan announces, sections 1 to 4: spectral
%     measure μ_f of a measurable function, the decreasing rearrangement φ_f
%     and formula (4) for it, the inequalities (5), the functions Φ_f, Ψ_f,
%     Δ_f and the determinant Δ(f) with Δ(fg) = Δ(f)Δ(g), Proposition 2
%     (∫ fg ≤ ∫ φ_f φ_g); then fonctions de Pólya (P1-P3, P'3) and their
%     Hahn-Banach representation P(f) = sup_i (a_i + ∫ φ_i φ_f), fonctions de
%     Weyl W(f) with W_α(f) = ∫ α(φ_f) and the N_p example; Prop 4 (∫_0^t f ≤
%     ∫_0^t g for all t implies P(f) ≤ P(g)) proved by parts, and Corollaire 1
%     (Δ_f ≤ Δ_g implies W(f) ≤ W(g)); then traces on a von Neumann algebra,
%     the ideal a of definition, its norm closure b, and Proposition 1 on
%     when the projectors and the K(M) of a commutative subalgebra fall in a.
%     Pages 6-10 are in black ink and mostly readable; pages 11-20 are in blue
%     and the prose is largely lost — the apparatus says so rather than
%     inventing. Page 19 is cancelled by two long diagonals.
%
% A later hand in blue ballpoint annotates pages 6, 8-10 (circling
% "Proposition 1", adding exp Ψ_{|f|}, "sur ]0,|M|[", "(|M| < +∞)"), and a
% black-ink hand corrects Ω to M throughout pages 19-20. Both are recorded
% once and not attributed. His own underlinings are set with \emph.
%
% The dating is the inventory's own for the shelfmark. Its group,
% « MATHÉMATIQUES / RECHERCHE / Travaux », carries no date.
\documentclass[11pt,a4paper]{article}
% Le préambule commun à toutes les transcriptions du fonds Grothendieck.
%
% Il tient en une idée : **l'incertitude se compose, elle ne se résout pas.**
% Une transcription qui lisse un mot douteux a détruit la seule chose pour
% laquelle on la faisait — un lecteur doit pouvoir savoir, à chaque endroit,
% ce qui était sur le feuillet et ce qui a été ajouté. Les six macros qui
% suivent sont donc l'appareil critique, et non de la décoration.
%
% Les mêmes commandes sont reconnues par `scripts/render.mjs`, qui produit la
% vue de lecture du site. Une macro ajoutée ici doit l'être aussi là-bas,
% faute de quoi le rendu échouera bruyamment — ce qui est voulu : un rendu
% silencieusement faux est pire qu'un rendu absent.

\ProvidesPackage{grothendieck}[2026/08/08 Transcriptions du fonds Grothendieck]

\usepackage{amsmath,amssymb,amsthm}
\usepackage{mathrsfs}
\usepackage{stmaryrd}
% Les diagrammes commutatifs sont partout dans ce fonds : c'est la langue
% même de la géométrie algébrique de Grothendieck, et une transcription qui
% les rendrait en prose ne transcrirait rien.
\usepackage{tikz-cd}
% Français par défaut — c'est la langue des feuillets — mais l'anglais est
% chargé aussi : la traduction et le résumé sont des documents anglais, et la
% typographie française (espaces avant les deux-points) y serait une faute.
% Ces documents basculent par \selectlanguage{english} en tête de corps.
\usepackage[english,french]{babel}
\usepackage{fontspec}
\usepackage{geometry}
\geometry{a4paper,margin=25mm,marginparwidth=28mm}
\usepackage{marginnote}
\usepackage{xcolor}
% ulem, not soul. `soul`'s \st and \ul are built on a character-by-character
% scan that XeLaTeX's font machinery breaks: the document fails with an
% undefined control sequence rather than degrading. ulem does the same two jobs
% and is Unicode-engine safe. `normalem` stops it hijacking \emph, which the
% transcriptions use for Grothendieck's own emphasis.
\usepackage[normalem]{ulem}
\usepackage{hyperref}
\hypersetup{colorlinks=true,linkcolor=black,urlcolor=black,pdfborder={0 0 0}}

% --- Métadonnées du lot -----------------------------------------------------
% Reprises telles quelles par le script de rendu, qui les affiche en tête de
% la vue de lecture. Elles fixent ce que le fichier prétend couvrir : sans
% elles, une transcription flotte sans point d'attache dans un fonds de seize
% mille feuillets.

\newcommand{\folder}[1]{\gdef\grothFolder{#1}}
\newcommand{\batch}[1]{\gdef\grothBatch{#1}}
\newcommand{\leaves}[2]{\gdef\grothFirst{#1}\gdef\grothLast{#2}}
\newcommand{\foldertitle}[1]{\gdef\grothFoldertitle{#1}}
\gdef\grothFolder{?}\gdef\grothBatch{?}\gdef\grothFirst{?}\gdef\grothLast{?}\gdef\grothFoldertitle{}

% --- Le résumé --------------------------------------------------------------
%
% Il ouvre la lecture modernisée, sous le titre « Résumé », et il est composé
% en petit. Ce n'est pas un ornement : c'est le seul passage du document qui
% ne soit pas des mathématiques, et un lecteur doit pouvoir le reconnaître
% avant de l'avoir lu — puis le sauter s'il connaît déjà le sujet, ou s'y
% arrêter s'il ne le connaît pas. Le filet qui le suit marque la reprise du
% corps.

\newenvironment{resume}
  {\small\setlength{\parskip}{0.45em}}
  {\par\medskip\hrule\medskip}

% --- Mots-clés ---------------------------------------------------------------
%
% En anglais, en clôture du résumé de la lecture modernisée : le vocabulaire
% moderne sous lequel chercher ce que les feuillets construisent. Le manifeste
% du site (`npm run manifest`) les extrait pour étiqueter la cote entière —
% cette ligne est la source unique des tags, il n'y a pas de fichier à côté.
\newcommand{\keywords}[1]{\par\smallskip\noindent
  {\footnotesize\textsc{Keywords} --- \textit{#1}}\par}

% --- L'appareil critique ----------------------------------------------------

% Le numéro de feuillet, en marge. C'est l'ancre qui permet de régler un
% désaccord sur une formule en regardant *un* feuillet plutôt qu'une chemise
% de six cents. Le site s'en sert aussi pour tourner les pages du fac-similé
% quand on fait défiler la transcription.
\newcommand{\leaf}[1]{%
  \marginnote{\footnotesize\color{gray!70}\textsf{#1}}%
  \label{leaf:#1}\ignorespaces}

% Illisible : on ne devine pas. Un mot inventé qui se lit comme les autres est
% le pire résultat possible.
\newcommand{\ill}{\textcolor{red!60!black}{[\,\ldots\,]}}

% Lecture douteuse : le mot est donné, et signalé comme incertain.
\newcommand{\uncertain}[1]{\uline{#1}}

% Ajout de l'éditeur — une lettre manquante, un mot sous-entendu.
\newcommand{\add}[1]{\textcolor{blue!50!black}{[#1]}}

% Biffé par Grothendieck lui-même. Ce qu'il a rayé fait partie du document :
% une rature dit souvent où la pensée a changé de direction.
\newcommand{\struck}[1]{\sout{#1}}

% Note du transcripteur, sur l'état matériel du feuillet.
\newcommand{\note}[1]{\par\smallskip{\small\color{gray!60!black}\textit{[#1]}}\par\smallskip}

% Note marginale de Grothendieck, distincte du corps du texte.
\newcommand{\marginal}[1]{\par\smallskip\noindent\hspace{1em}%
  {\small\color{gray!60!black}#1}\par\smallskip}

% --- Titre ------------------------------------------------------------------

\AtBeginDocument{%
  \begin{center}
    {\large\bfseries Fonds Alexandre Grothendieck --- cote n\textsuperscript{o}~\grothFolder}\\[2pt]
    {\itshape\grothFoldertitle}\\[4pt]
    {\small feuillets \grothFirst--\grothLast\ (lot \grothBatch)}\\[2pt]
    {\footnotesize\color{gray!60!black}Transcription établie d'après le fac-similé numérisé
     par l'Université de Montpellier.\\
     Les lectures incertaines sont soulignées, les illisibles notées [\,\ldots\,],
     les ajouts entre crochets.}
  \end{center}
  \vspace{1em}}

\endinput


\folder{1}
\batch{1}
\pages{1}{20}
\dating{1953}
\watermark{Édition de démonstration}
\foldertitle{Analyse fonctionnelle : notes manuscrites (s.d.), lettre (1953)}

\begin{document}

\section{Plan des notes}
\note{Titre de l'éditeur. La page 2 est le plan, de sa main, des sept numéros dont les pages 6 à 20 donnent les quatre premiers.}

\page{2}
\note{Le feuillet, plié en deux, porte dans sa moitié supérieure le plan des notes, de sa main : sept numéros, deux fois renumérotés — le « N° 5 » primitif est biffé, le « 6 » suivant corrigé en 5, le « 7 » final surchargé.}

\begin{itemize}
  \item N° 1 --- Notions sur les réarrangements de fonctions
  \item N° 2 --- Fonctions de Weyl.
  \item N° 3 --- Inégalités \struck{de convexité \ill{} liées aux réarrangements} \note{au-dessus du mot biffé « convexité » est écrit « élémentaires », au-dessous « diverses » ; et, sous la ligne, « $(A \in \bar{a})$ », repris à droite « $(A \in \tilde{\bar{a}})$ » et relié d'un trait à la ligne du N° 4}
  \item N° 4 --- Mesures spectrales, réarrangement \struck{spectral} d'un opérateur.
  \item \struck{N° 5} --- \struck{Premières propriétés des réarr\add{angements} \ill{} d'un \ill{}}
  \item N° 5 \note{écrit sur un « 6 »} --- Compléments sur le déterminant
  \item N° 6 --- Les inégalités \struck{de convexité dans les anneaux d'opérateurs} \note{au-dessus du passage biffé : « \uncertain{fondamentales} \struck{pour} $AB$ »}
  \item N° 7 --- Autres inégalités.
\end{itemize}

\section{Endomorphismes d'une $C^{*}$-algèbre permutables aux translations}
\note{Titre de l'éditeur : le feuillet n'en porte aucun. Les pages 3 à 5, à l'encre bleue sur feuillets pliés, forment un fragment distinct des notes sur les réarrangements : $C$ est une $C^{*}$-algèbre, $C''$ son bidual, $R$ une algèbre de von Neumann où $C$ est plongée, et le fragment identifie l'algèbre $\Gamma$ des $u \in C''$ tels que $uC \subset C$, $Cu \subset C$ à l'algèbre $\Gamma_1$ des $x \in R$ de même propriété. La main est rapide ; les formules passent, la prose entre elles souvent non.}

\page{3}
\note{Les feuillets portent, au crayon de l'archiviste, une numérotation qui court de deux en avance sur la position dans la cote à partir d'ici : ce feuillet porte 5, le suivant 6, et ainsi de suite jusqu'à 22 pour le dernier du lot. Les numéros employés pour \texttt{page} sont ceux de la position, que compte l'inventaire (149 pages) et que suit le fac-similé.}

Soit $C$ une $C^{*}$-algèbre, \ill{} \uncertain{déterminer} les \uncertain{endomorphismes} linéaires continus $T$ \ill{} \uncertain{aux translations à gauche}. Soit $T$ une telle, son \uncertain{bitransposé} $T''$ est un endomorphisme \add{de $C''$} \ill{} \uncertain{permutable} \ill{} \uncertain{aux translations} $x \to xy$ \struck{($y \in C$)} \ill{} pour $y \in C$, donc pour \struck{tout} $y \in C''$ ($T''(xy) = T''(x)\,y$) \struck{\ill{}} \ill{} \note{une ligne interpolée au-dessus et à droite, non lue}. Donc de la forme $x \to ux$ ($u \in C''$), \ill{}. \ill{} : l'ens. $\Gamma_g$ des $u$ de $C''$ tels que $uC \subset C$. On voit donc que l'ensemble d'endomorphismes considéré est isomorphe, avec sa structure d'algèbre normée, à $\Gamma_g$. \struck{\ill{}} \uncertain{« hermitien »} signifie $x^{*}T(x)$ hermitien pour tout $x \in C$ ; \uncertain{« positif »} signifie $x^{*}T(x) \geq 0$ pour \struck{tout} $x \in C$. Soit \struck{\ill{}} $\Gamma$ l'ens. des $u \in C''$ tels que $u\underline{C} \subset \underline{C}$, $\underline{C}u \subset \underline{C}$ \note{le $C$ est ici souligné, à quatre reprises}. C'est une sous-$C^{*}$-algèbre \uncertain{involutive}. \ill{} \uncertain{hermitien} à $\Gamma_g$ \ill{} dans $\Gamma$.

Supposons maintenant $C$ plongé dans la $C^{*}$-algèbre $R$, et soit $\Gamma_1$ l'ens. des $x \in R$ tels que $xC \subset C$, $Cx \subset C$, c'est une sous-$C^{*}$-algèbre de $R$. Pour tout $u \in \Gamma_1$, soit $\varphi(u)$ l'élément de $C''$ correspondant : l'\uncertain{endomorphisme} $T_u x = ux$ dans $C$. Si $u$ hermitien, $\varphi(u)$ est hermitien, donc dans $\Gamma$, donc pour \struck{tout} $u \in \Gamma_1$, $\varphi(u) \in \Gamma$. De plus, $\varphi$ est évidemment une représentation d'algèbres, compatible avec l'involution (car \ill{}

\page{4}
transforme hermitiens en hermitiens). \uncertain{Cela} \uncertain{donne} \ill{} \uncertain{vérification} un \emph{\uncertain{isomorphisme}} entre $\Gamma$ et $C'$, \uncertain{partiel} \ill{}. Car on a
\[
  \varphi(a) = \lim_{x \to 1} ax = \lim_{x \to 1} xa,
\]
\note{sous chaque $\lim$, en trois lignes : « $x \to 1$ », « \uncertain{filt} », « $C$ » — le filtre des unités approchées de $C$, nommé ainsi plus bas ; sous le premier, une première version est raturée}
\uncertain{d'où l'on tire} \ill{} \uncertain{donc} (") d'où
\[
  \varphi\langle a, f\rangle = \lim_{x \to 1} \langle ax, f\rangle
\]
($a \in \Gamma_1$, $f \in C'$), et \ill{} \uncertain{compte tenu} \ill{} \uncertain{compte tenu de} \ill{}, \uncertain{involution} \ill{}, et \uncertain{prolonge} \ill{} \uncertain{donné} \struck{\ill{}} sur $\mathbf{C} \times C'$.

Supposons maintenant \ill{} \uncertain{de} \uncertain{Neumann}, \note{ajouté au-dessus : « et $C$ dense dans $R$ »} d'où \uncertain{une} \uncertain{application} \uncertain{canonique} $\psi : C'' \to R$ (\uncertain{en identifiant} $C''$ \ill{} l'adhérence \uncertain{faible} de $C$ dans $R$). Évidemment $\psi(\Gamma) \subset \Gamma_1$. D'autre part, \uncertain{prouvons} $\varphi\psi\struck{(u)} = 1$ sur $\Gamma$, car soit $u \in \Gamma$, \struck{\ill{}} et $v = \psi(u) \in \Gamma_1$, \ill{} pour \struck{tout} $f \in R_{*}$
\[
  \langle v, f\rangle = \langle u, f_1\rangle
\]
($f_1$ \uncertain{image} de $f$ dans $C'$ ; \ill{} \uncertain{on identifie} $R_{*} \subset C'$). J'affirme que $\langle v, f\rangle = \langle \varphi(v), f_1\rangle$ [\ill{} \uncertain{ceci} \uncertain{prouve} \struck{$\varphi\psi$} $\varphi(v)$ \ill{} $u$ \uncertain{orthogonal} \ill{} : $R_{*}$] \uncertain{cela} se voit \uncertain{facilement} sur $\langle \varphi(v), f_1\rangle = \lim_{x \to 1} \langle vx, f\rangle$ on \uncertain{prend} ici le filtre \uncertain{des unités} \uncertain{approchées} \ill{} \uncertain{le cas de} $\mathbf{C}$. On est ramené à \uncertain{prouver} : un élément de $\underline{\Gamma}$ orthogonal à $R_{*}$ \struck{est nul} (i.e. dans l'image dans $R$ est $0$) est nul. C'est \uncertain{immédiat}

\page{5}
car si $x$ est un tel élément, alors pour $y \in C$ \ill{} $\varphi(xy) = \varphi(x)y = 0$, d'où (comme $xy \in C$, \uncertain{d'où} $\varphi(xy) = xy$) $xy = 0$, et ceci pour tout $y \in C$, d'où $x = 0$. --- \emph{Donc $\varphi$ est un isomorphisme de $\Gamma$ dans $\Gamma_1$}, \struck{\ill{}} \emph{et $\psi$ l'inverse} \note{« et $\psi$ l'inverse » est écrit sur le passage biffé}. \ill{} Plus précisément, prouvons que $\psi\varphi = 1$, soit donc $u \in \Gamma_1$, $\varphi u = v \in C''$ \uncertain{est} \uncertain{donné} par $\langle v, f\rangle = \langle u, f\rangle$ ($f \in C'$) \uncertain{tandis que} $\psi v$ est \uncertain{donné} par $\langle \psi v, f\rangle = \langle v, f_1\rangle$ ($f \in R_{*}$) i.e. $\langle \psi v, f\rangle = \langle u, f_1\rangle$ ($f \in R_{*}$), \uncertain{d'autre part} \ill{} \ill{} $\langle u, f_1\rangle = \langle u, f\rangle$, d'où $\langle \psi v, f\rangle = \langle u, f\rangle$, i.e. $\psi v = u$. CQFD.

\section{Mesures spectrales et réarrangements de fonctions}
\note{Titre de sa main en tête de la page 6 ; c'est le « N° 1 » du plan de la page 2. Une main postérieure a entouré au stylo bille bleu le numéro de la proposition 1.}

\page{6}
\marginal{$|M|$, $L^{1}(M)$, $L^{\infty}(M)$ ; $K(M)$, \uncertain{$K^{\infty}(M)$}, $C_0(M)$ ; \struck{\ill{}}}

Soit $f$ une fonction mesurable complexe sur l'espace mesuré $M$ (muni de la mesure $m$). On dit que $f$ \emph{admet une mesure spectrale} si l'image de $m$ par $f$ \uncertain{sur} \ill{} existe. Cela signifie donc aussi que pour $0 < a < b < +\infty$, $|f|^{-1}([a,b])$ est intégrable (donc cette condition ne dépend que de $|f|$). \uncertain{Cette} \uncertain{mesure} \uncertain{sur} $\underline{\mathbf{C}^{*}}$ \uncertain{est} \uncertain{appelée} \emph{mesure spectrale de $f$}, et notée $\mu_f$ :
\begin{align*}
  \text{(1)}\qquad & \mu_f = f(m) \quad\text{i.e.} \\
  \text{(2)}\qquad & \langle \alpha, \mu_f\rangle = \langle \alpha\circ f, m\rangle = \int \alpha\circ f\, dm
\end{align*}
(\struck{$\alpha$ fonction} \ill{} \uncertain{fonction} continue à support compact sur $\underline{\mathbf{C}^{*}}$). $\mu_f$ est \emph{positive}. Elle est portée par $\underline{\mathbf{R}_{*}}$ (resp. par $\mathbf{R}^{+}_{*}$) si et seulement si $f$ est réelle (resp. positive). \ill{} une fonction \uncertain{sur} \ill{} espace mesuré $M$ et une fonction \uncertain{$g$} sur l'espace mesuré $N$ \uncertain{sont} dites \emph{spectralement} \emph{équivalentes} si elles \uncertain{ont} \uncertain{même} mesure spectrale.

\emph{Exemples}. Si l'ensemble \uncertain{$f^{-1}(\ldots)$} \ill{} est intégrable, alors $\mu_f$ existe et est bornée. C'est \uncertain{en particulier} le cas si $f$ est à support intégrable \ill{} $m$ est bornée \ill{} $f \in C_0(M)$, \ill{} \note{deux lignes interpolées ici, à demi effacées : « \ldots\ pour tout $\varepsilon$, $|f|^{-1}(]\varepsilon, +\infty[)$ est intégrable »}.

\emph{Proposition 1} Une fonction positive décroissante $\varphi$ sur $\mathbf{R}^{+}_{*}$ a une mesure spectrale si et seulement si elle est nulle à l'infini ; \uncertain{c'est-à-dire} $\mu_\varphi$ est une mesure positive sur $\mathbf{R}^{+}_{*}$ telle que $\int_1^\infty d\mu_\varphi < +\infty$. Réciproquement pour toute mesure positive $\mu$ sur $\mathbf{R}^{+}_{*}$ telle que $\int_1^\infty d\mu < +\infty$, existe une fonction décroissante \add{positive} $\varphi$ \note{« positive » ajouté au-dessus de la ligne} et une seule (à une \emph{fonction} \struck{\ill{}} \uncertain{négligeable} près) qui \uncertain{admet} $\mu$ \uncertain{pour} \uncertain{mesure} spectrale.

On a donc pour presque tout $t$
\[
  \text{(3)}\qquad t = \int_s^\infty d\mu(s) \qquad (s = \varphi(t))\,.
\]

\page{7}
les $t$ exceptionnels étant ceux pour lesquels \uncertain{la} \ill{} de \ill{} \uncertain{de} (1) est ambiguë, i.e. qui \uncertain{sont} \ill{} par \uncertain{la} \ill{} \uncertain{telle} \ill{}, \uncertain{ceux} \ill{} \uncertain{tel} $t$. Ainsi $\varphi$ \uncertain{apparaît} \uncertain{comme} \uncertain{une} fonction \uncertain{inverse} de
\[
  s \to \int_s^\infty d\mu(s)\,.
\]

\struck{Définition}

\emph{Corollaire} Soit $f$ une fonction mesurable positive sur l'espace mesuré $M$. Pour que $f$ soit spectralement équivalente à une fonction décroissante \uncertain{$\varphi$} sur $\mathbf{R}^{+}_{*}$, il faut et il suffit que pour tout $a > 0$, $f^{-1}([a, +\infty[)$ soit intégrable. Alors $\varphi_f$ \note{le $\varphi$ porte un indice raturé} est \uncertain{unique} \uncertain{et} \uncertain{déterminée} \ill{} \struck{Définition} On l'appelle le \emph{réarrangement décroissant de $f$}, noté $\varphi_f$. Si \struck{\ill{}} $g$ est une autre fonction mesurable positive, sur un \uncertain{autre} espace mesuré $N$, satisfaisant aux conditions ci-dessus, alors $f$ et $g$ sont spectralement \struck{\ill{}} équivalentes si et seulement si elles ont même réarrangement décroissant. \struck{\ill{}} \note{une ligne biffée puis cernée d'un trait ; au-dessous, « $\varphi_f(t)$ est \uncertain{déterminé} \ill{} » également biffé}
\[
  \text{(4)}\qquad \varphi_f^{-}(t) = \inf_{\substack{E \subset M\\ m(E) \leq t}} \|(1-E)f\|_\infty \qquad
  \varphi_f^{+}(t) = \inf_{\substack{E \subset M\\ m(E) < t}} \|(1-E)f\|_\infty
\]
\note{le « $\leq$ » du premier infimum est surchargé ; sous le second, il est écrit « $-(E) < t$ » — sans doute un $m$ mal formé}
(\struck{$\varphi_f^{-}(t)$} \ill{} $\varphi^{-}(t) = \lim_{s \to t,\ s < t} \varphi(s)$, $\varphi_f^{+}(t) = \lim_{s \to t,\ s > t} \varphi(s)$ \note{la seconde formule est écrite au-dessus d'une première version biffée}) \uncertain{où} \uncertain{on identifie} $E$ avec sa fonction caractéristique. \ill{} \uncertain{à} \ill{} : (4). De (4) on tire les \struck{\ill{}} \uncertain{propriétés} \uncertain{suivantes} (pour des fonctions mesurables positives \ill{} \struck{\ill{}} \note{deux lignes biffées, la seconde cernée d'un long trait ; en tête d'une troisième, « \uncertain{mesuré} $M$ », puis l'accolade qui suit}) :
\[
  \text{(5)}\qquad \left\{
  \begin{array}{ll}
    \varphi_f \leq \varphi_g & \text{si } f \leq g \\
    \varphi_{\lambda f} = \lambda\varphi_f & \text{si } \lambda \geq 0 \\
    \varphi_{f^\alpha} = (\varphi_f)^\alpha & \text{si } \alpha \geq 0 \\
    \varphi_{f+g} \leq \varphi_f + \varphi_g & \\
    \varphi_{fg} \leq \varphi_f\varphi_g &
  \end{array}\right.
\]
\ill{} \uncertain{la} \uncertain{première} \ill{}, (si $\varphi_g$ existe et $f \leq g$, $\varphi_f$ existe \ill{} ; $\varphi_{f^\alpha}$, $\varphi_{f+g}$ et $\varphi_{fg}$ existent). \note{« et $f \leq g$, $\varphi_f$ » est ajouté au-dessus de la ligne, et $\varphi_{f^\alpha}$, $\varphi_{f+g}$ en surcharge}

\marginal{dans une bulle en marge gauche, à hauteur de (4) : $\int f\,dm = \int s\,d\mu(s)$, d'où $\int f\,dm = \int \varphi_f\,dt$ ; $N_p(f) = N_p(\varphi_f)$, etc.}

\marginal{dans une seconde bulle, reliée d'une flèche à la ligne $\varphi_{f^\alpha} = (\varphi_f)^\alpha$ : \ill{} plus \uncertain{généralement}, \uncertain{pour} \uncertain{$\alpha$} \uncertain{croissante} \ill{} $\mathbf{R}_{*}^{+}$ \struck{\ill{}} $\varphi_{\alpha(f)} = \alpha(\varphi_f)$ (\ill{}) \ill{} ; puis, biffées au stylo bleu, deux lignes dont on lit « $\varphi_{|f|} = \ldots \varphi_f \ldots$ »}

Soit $\varphi$ une fonction décroissante sur $\mathbf{R}^{+}_{*}$, admettant la mesure spectrale $\mu_\varphi$. Pour que

\page{8}
$\varphi$ soit \uncertain{sommable} au voisinage de $0$, il faut et il suffit \uncertain{qu'il} \ill{} et \uncertain{sommable} \ill{} $\int_1^\infty d\mu$ \ill{} $\int_0^\infty d\mu(s) < +\infty$, i.e.
\[
  \int_1^\infty (s-1)\, d\mu(s) < +\infty, \qquad\text{ou encore}\qquad
  \text{(6)}\quad \int^\infty s\, d\mu(s) < \struck{\ill{}} +\infty
\]
\note{la formule (6) est encadrée} \uncertain{condition} \uncertain{vérifiée} \uncertain{quand} \ill{} \struck{\ill{} Si $\mu = \mu_f$ \ill{} $f$ \ill{} \uncertain{sommable}, i.e. \ill{} l'espace mesuré $M$} quand $\mu = \mu_f$, $f$ étant une \uncertain{fonction} \ill{} \uncertain{sommable} \ill{} \uncertain{sur} \ill{} \struck{Si \uncertain{alors} $\varphi_f$ \uncertain{satisfait} \ill{}} Si $\varphi_f$ est sommable au voisinage de $0$, on pose
\[
  \text{(7)}\qquad \Phi_f(t) = \int_0^t \varphi_f(s)\, ds
\]
$\Phi_f(t)$ est une fonction continue concave \uncertain{croissante} de $t$. Considérons $\log\varphi_f(s)$, c'est une fonction \struck{décroissante} \add{\uncertain{continue}} \note{« continue » écrit au-dessus du mot biffé}, à valeurs dans $[-\infty, +\infty[$. \uncertain{Supposons} \ill{} $\log\varphi_f(t)$ \uncertain{sommable} \uncertain{au} voisinage de \ill{}, i.e. $\log f$ sommable dans $f^{-1}([1, +\infty[)$, \struck{$\log\varphi_f(t)$ \ill{}} \note{trois lignes biffées, la dernière cernée d'un trait, où l'on lit « $\int \varphi_f(s)$ \ldots » et « $\log\varphi_f(t)$ »} (à fortiori $\varphi_f$ \ill{}) \ill{}
\[
  \text{(8)}\qquad \Psi_f(t) = \int_0^t \log\varphi_f(s)\, ds
\]
\note{le signe $=$ est surchargé ; au-dessus, un mot raturé} Si $\varphi_f$ ne \uncertain{s'annule} \uncertain{pas} (i.e. si $f^{-1}(\{0\})$ \uncertain{n'est} \uncertain{pas} \ill{} \note{le passage biffé d'un trait épais porte « $f^{-1}(]0, +\infty[)$ »}) \uncertain{fonction} \uncertain{positive} \uncertain{finie} \ill{} \uncertain{intégrable}), $\Psi_f(t)$ est \uncertain{une} \uncertain{fonction} \ill{} \uncertain{concave}. \uncertain{Sinon}, si $t_0 = |f^{-1}(\complement\{0\})|$ \note{le complémentaire est noté d'un $C$ barré} \uncertain{du} \uncertain{moins} $\varphi_f$ \ill{}, \uncertain{c'est} \uncertain{le} \uncertain{point} \uncertain{à} partir \uncertain{duquel} $\Psi_f(t)$ est égal à $-\infty$. Enfin, on pose
\[
  \text{(9)}\qquad \Delta_f(t) = \exp\Psi_f(t)
\]
\note{à droite de (9), au stylo bleu et d'une autre main : « $\exp\Psi_{|f|}(t)$ » ; le même stylo retouche une lettre deux lignes plus bas} Sous les conditions \uncertain{envisagées} sur $f$, $\Delta_f(t)$ est une fonction continue (\struck{positive}) \uncertain{croissante} \note{« croissante » écrit au-dessus} positive de $t$ ; elle ne \uncertain{peut} s'annuler que si $t_0 = |f^{-1}(\complement\{0\})|$ est fini, \uncertain{alors} $\Delta_f(t)$ \uncertain{s'annule} à partir de $t_0$.

Les formules (5) donnent des inégalités correspondantes \uncertain{pour} \struck{pour} $\Phi_f$, $\Psi_f$, $\Delta_f$. \uncertain{Les} \uncertain{plus} \uncertain{importantes} \uncertain{sont}
\begin{align*}
  \text{(10)}\qquad & \Psi_{fg} \leq \Psi_f + \Psi_g \qquad\text{i.e.} \\
  \text{(11)}\qquad & \Delta_{fg} \leq \Delta_f\,\Delta_g
\end{align*}

\page{9}
\struck{Supposons de nouveau \ill{} $\varphi_f$ existe}. On pose, pour une fonction mesurable \struck{positive}
\[
  \text{(12)}\qquad \Delta(f) = \Delta_{|f|}(|M|) = \struck{\exp} \int_M \log|f|\, dm
\]
\note{l'indice $|f|$ de $\Delta$ est repassé au stylo bleu par la main postérieure} \struck{$\Delta$ existe pour} \struck{donc} $\Delta(f)$ existe \struck{si $\Delta_f$ existe} \uncertain{ssi} $\log|f|$ est sommable sur l'ens. des points où il est positif. $0 \leq \Delta(f) < +\infty$. $\Delta$ jouit des propriétés d'un (\add{module de} \note{« module d\ldots » ajouté au-dessus}) déterminant, en particulier
\[
  \text{(13)}\qquad \Delta(fg) = \Delta(f)\,\Delta(g) \qquad \Delta(1) = 1
\]
\note{un trait ondulé sépare ici le texte}

Supposons que $|M| < +\infty$. On voit \uncertain{aussitôt} que si $f$ est une fonction mesurable positive sur $M$, \struck{admettant un réarrangement \uncertain{décroissant} $\varphi_f$, alors} \struck{il en est de même} \struck{$\varphi_{1+f}(t)$}
\[
  \text{(14)}\qquad \varphi_{1+f}(t) = 1 + \varphi_f^{-}(t) \qquad \text{si } 0 < t \leq |M|
\]
\note{le « $1$ » de $1+f$ et de $1+\varphi_f^{-}$ est penché et se lit aussi $\lambda$ ; on garde $1$, que (16) écrit sans ambiguïté} (si $\varphi_f$ existe ; \uncertain{si} \uncertain{non} $|M| = \infty$, $\varphi_{1+f}$ n'existe pas \ill{} $f \geq 0$ \ill{} \struck{$1 + \varphi_f(t)$ existe, donc} \ill{} \uncertain{donc} (14)). Alors $\log(\varphi_{1+f}) = \log(1 + \varphi_f)$, \struck{donc $\log(1+\varphi_f)$ \ill{}} \note{cerné d'un trait} \uncertain{fonction} positive. Si on \uncertain{suppose} que $\Phi_f$ existe, alors $\log(1+\varphi_f)$ est sommable au voisinage de $0$ (et donc \ill{} \uncertain{partout}). \uncertain{Ce} \uncertain{qui} \uncertain{permet} \uncertain{de} \uncertain{définir}
\[
  \text{(15)}\qquad \Delta_{1+f}(t) = \exp\int_0^t \log(1 + \varphi_f(s))\, ds \qquad 0 \leq t \leq |M|
\]
cette notation est \uncertain{cohérente} \uncertain{avec} \uncertain{la} \uncertain{définition} \ill{} \uncertain{Car} \ill{} \uncertain{d'ailleurs}
\[
  \text{(16)}\qquad \Delta_{1+f}(|M|) = \Delta(1+f)
\]
\note{un second trait ondulé, et une accolade en marge droite qui embrasse le passage de (14) à (16)}

Soit $\mu$ une mesure positive sur $\mathbf{R}_{*}$, telle que $\int_a^\infty d\mu < +\infty$ (\ill{} \uncertain{ne} \uncertain{dépendant} \ill{}) \note{en interligne, au stylo bleu : « \struck{$|M| \geq \|\mu\|$} Supposons $\|\mu\| = \ldots$ \uncertain{ou} $0 \leq \|\mu\| \leq |M| < +\infty$ » ; en marge gauche, dans un ovale biffé, de la même main bleue : « Le \uncertain{mieux} : \ldots\ $\|\mu\| < +\infty$ ! »} (Alors il existe une \uncertain{fonction} \uncertain{décroissante} \ill{} $\mathbf{R}_{*}^{+}$ (\struck{\uncertain{resp.} $]0, |M|[$} \uncertain{sur} $]0, |M|[$) \note{parenthèse ajoutée au stylo bleu}) \ill{} \uncertain{fonction} décroissante $\varphi$ et une seule sur $\mathbf{R}^{+}_{*}$ telle que $\mu_\varphi = \mu$. Cette $\varphi$ est \uncertain{une} \uncertain{fonction} \uncertain{positive} si et seulement si $\mu$ est portée par $\mathbf{R}^{+}_{*}$. Soit maintenant $f$ telle

\page{10}
fonction \emph{réelle} mesurable sur $M$ telle \struck{\ill{}} (si $|M| < +\infty$) \uncertain{telle} \struck{que $\varphi_f$ soit \ill{} \uncertain{partie} \uncertain{positive} \ill{}, il suffit que \ill{} pour $\varphi_g$ \ill{} bornée, donc que \ill{} $f$ \ill{} que} \struck{$|f^{-1}(\complement\{0\})| < \infty$} \note{trois lignes biffées, la dernière au stylo bleu}. Alors il existe une fonction décroissante $\varphi_f$ et une seule \add{sur $]0, |M|[$} \note{ajouté au stylo bleu, dans l'interligne} \uncertain{spectralement} \uncertain{équivalente} \uncertain{à} $f$. On l'appelle encore \uncertain{le} \uncertain{réarrangement} \uncertain{décroissant} de la fonction réelle $f$. Le \emph{résultat suivant est classique} :

\emph{Proposition 2} Soient $f, g$ deux fonctions réelles mesurables sur $M$, \struck{telles} ($|M| < +\infty$) \note{ajouté au stylo bleu} \uncertain{telles} \uncertain{que} $fg$ soit sommable, il \uncertain{faut} \uncertain{que} $\varphi_f\varphi_g$ le soit, et \ill{}. alors
\[
  \text{(17)}\qquad \int fg\, dm \leq \int_0^{|M|} \varphi_f\varphi_g\, ds
\]
\note{sous (17), un mot au stylo bleu, non lu ; le reste de la page est blanc}

\section{Fonctions de Weyl}
\note{Titre de sa main : « 2. Fonctions de Weyl », page 11. De la page 11 à la page 17, l'encre est bleue ; le crayon des marges est d'une autre campagne.}

\page{11}
\marginal{au crayon, en marge gauche du titre : $K^\infty$, $L^\infty(K)$ ; $\|M\|$ \ill{} ; $K^\infty_k$ ; \ill{}}

\emph{Définition 1} Soit \struck{$M$ un espace mesuré} $I = ]0, \omega[$ ($0 < \omega \leq +\infty$) \note{écrit au-dessus du passage biffé} \uncertain{On} \uncertain{appelle} \emph{fonction de Pólya} sur \struck{$M$} \add{$I$} \struck{une fonction $P$ réelle} \uncertain{une} \uncertain{fonction} $P$ réelle définie sur l'espace $K^\infty(\bar{I})$ des fonctions mesurables \uncertain{à} \uncertain{support} \uncertain{compact}, \uncertain{ayant} les propriétés suivantes :
\begin{itemize}
  \item[$P_1$)] $P$ est \emph{symétrique}, i.e. $P(f) = P(g)$ si $f$ et $g$ sont \uncertain{spectralement} \uncertain{équivalentes}
  \item[$P_2$)] $P$ est convexe, i.e. $P(\lambda f + (1-\lambda) g) \leq \lambda P(f) + (1-\lambda) P(g)$ si $0 \leq \lambda \leq 1$
  \item[$P_3$)] Pour tout compact $K$, \struck{$P$ est} \uncertain{et} \uncertain{toute} \uncertain{suite} $(f_i)$ \uncertain{dans} $K^{\infty}_{k}(K)$ \note{l'indice se lit $k$ ; la marge au crayon porte le même $K^\infty_k$} \uncertain{tendant} \uncertain{vers} \uncertain{une} \uncertain{limite} $f$ \uncertain{en} \uncertain{mesure}, on a $P(f) \leq \varliminf P(f_i)$
\end{itemize}
\note{une accolade à gauche embrasse $P_1$ à $P_3$}

En fait, $P_3$ est \uncertain{équivalente} \uncertain{à} \uncertain{la} \uncertain{condition} \uncertain{plus} \uncertain{forte} \ill{} \uncertain{à} \uncertain{savoir} :
\begin{itemize}
  \item[$P'_3$)] ($P$ est \emph{semi-continue} \emph{inférieurement} \uncertain{sur} \struck{tout compact $K$} \uncertain{pour} \uncertain{la} \uncertain{topologie} \uncertain{faible} \ill{} \uncertain{pour} \uncertain{tout} \uncertain{compact} $K$ \note{ajouté au-dessus}, \ill{} \uncertain{sur} \uncertain{les} \uncertain{parties} \uncertain{bornées} \uncertain{de} $L^\infty(K)$
\end{itemize}
Il faut montrer que pour \uncertain{tout} $a$, \struck{$P^{-1}([a, +\infty[)$} \uncertain{l'ensemble} $A$ \uncertain{des} \note{ajouté au-dessus} $f \in L^\infty(K)$ telles que $P(f) \leq a$ \uncertain{est} \struck{\ill{}} \uncertain{fermé} \uncertain{pour} \uncertain{la} \uncertain{topologie} \ill{} \struck{\ill{} \uncertain{Bourbaki} \ill{} \uncertain{Dieudonné} \ill{}} \note{quatre lignes biffées et cernées d'un trait, dont une reprend « $L^\infty(K)$ »} \ill{} \uncertain{ce} \uncertain{qui} \uncertain{résulte} \uncertain{de} $P_3$ \uncertain{il} \uncertain{est} \struck{\ill{}} \uncertain{fermé} \uncertain{pour} \uncertain{la} \uncertain{convergence} \uncertain{en} \uncertain{mesure} \uncertain{sur} \uncertain{les} \uncertain{boules} \uncertain{de} $L^\infty(K)$ \uncertain{pour} \ill{} \uncertain{mesure}. Or on \uncertain{sait} \uncertain{que} \uncertain{un} \uncertain{ensemble} \uncertain{convexe} \ill{} $A$ \uncertain{est} \uncertain{dit} \uncertain{fermé} \ill{}, \ill{} \uncertain{un} \uncertain{...} \uncertain{fermé} \uncertain{dans} \uncertain{tout} \ill{} $B$ \ill{} \uncertain{Banach}, \uncertain{si} $A$ \uncertain{est} \uncertain{borné} \ill{} \uncertain{fermé} \uncertain{dans} $L^\infty(K)$. \uncertain{Il} \uncertain{est} \uncertain{donc} \ill{} \uncertain{fermé} \uncertain{dans} $L^1(K)$, \struck{donc} \uncertain{de} \uncertain{plus} \uncertain{il} \uncertain{est} \uncertain{fermé} \uncertain{dans} $L^1(K)$ (\uncertain{pour} \uncertain{la} \uncertain{topologie} \uncertain{faible}, \uncertain{donc} \uncertain{pour} \ill{}

\page{12}
\uncertain{sur} $]0, +\infty[$, \uncertain{si} \uncertain{on} \uncertain{note} \ill{} \uncertain{une} fonction $P$ \uncertain{sur} $K^\infty(]0, +\infty[)$ \uncertain{est} \uncertain{une} \uncertain{fonction} \uncertain{de} \uncertain{Pólya} si et seulement si \uncertain{pour} \uncertain{tout} $0 < \omega < +\infty$, \uncertain{sa} \uncertain{restriction} \uncertain{à} $L^\infty(]0, \omega[)$ \uncertain{est} \uncertain{une} \uncertain{fonction} \uncertain{de} \uncertain{Pólya}. --- \struck{\ill{} Soit \ill{} $M$ \ill{} $|M| < \infty$ \ill{} \uncertain{fonction} \uncertain{de} \uncertain{Pólya} \ill{} $L^\infty(]0, \omega[)$ \ill{} \uncertain{restriction} \uncertain{d'une} \uncertain{fonction} \ill{} $(0 < \omega < +\infty)$ \ill{} $K^\infty(]0, +\infty[)$ \ill{} Pólya \uncertain{sur} \ill{} Soit $M$ \ill{}} \note{sept lignes biffées, cernées de traits qui les relient entre elles ; on y lit encore « $f \in K^\infty(M)$ »}

\marginal{en marge gauche, de haut en bas : Pólya : $P(f) = \int \alpha(f)\,dt$ \uncertain{où} $\alpha$ \uncertain{fonction} \uncertain{convexe} \ill{} ; \uncertain{exemples} \ill{} $x \to e^x$ \ill{} ; \ill{} $= 0$ \uncertain{si} $x \leq 0$, \ill{} \uncertain{si} $x \geq 0$ ; \ill{} \uncertain{exemples} \ill{} $\to +\infty$ \ill{}}

\emph{Remarque} \uncertain{Fonction} \uncertain{de} \uncertain{Pólya} \ill{} \uncertain{voir} \uncertain{définition}.

\emph{Définition 2} \emph{Fonction de Weyl} \uncertain{Une} \uncertain{fonction} \uncertain{réelle} \uncertain{définie} \uncertain{sur} \ill{} \uncertain{est} \uncertain{dite} \struck{\ill{}}
\begin{itemize}
  \item[($W_1$)] $P(\exp f)$ \uncertain{est} \uncertain{une} \uncertain{fonction} \uncertain{de} \uncertain{Pólya} \uncertain{sur} $L^\infty(]0, \omega[)$.
\end{itemize}
Une fonction \uncertain{sur} $(K^{\infty})^{+}(]0, +\infty[)$ \uncertain{est} \uncertain{dite} \uncertain{fonction} de Weyl, si \uncertain{ses} \uncertain{restrictions} \uncertain{à} \uncertain{tout} $(K^\infty)^{+}(]0, \omega[)$ ($0 < \omega < +\infty$) \uncertain{le} \uncertain{sont}. \note{une accolade à gauche embrasse la définition}

\emph{Définition} \ill{} \emph{Weyl} \uncertain{sur} $L^{\infty +}(]0, \omega[)$ \ill{} \uncertain{satisfaisant} \uncertain{à} $P_1$, $P_3$, \struck{\ill{}} \add{\uncertain{continue} \uncertain{pour} \ill{}} \uncertain{et} \ill{} \uncertain{une} \uncertain{fonction} \uncertain{croissante} \uncertain{sur} \ill{}. \uncertain{Une} \uncertain{fonction} \uncertain{sur} $(K^\infty)^{+}(]0, +\infty[)$ \uncertain{est} \uncertain{une} \uncertain{fonction} \uncertain{de} \uncertain{Weyl} \uncertain{si} \uncertain{ses} \uncertain{restrictions} à \uncertain{tout} $(K^\infty)^{+}(]0, \omega[)$ \uncertain{le} \uncertain{sont}.

\ill{} \uncertain{les} \uncertain{fonctions} \uncertain{de} \uncertain{Weyl} \uncertain{sur} $L^{\infty}(]0, \omega[)$ \uncertain{et} \uncertain{les} \struck{\ill{}} \uncertain{fonctions} \uncertain{de} \uncertain{Pólya} \uncertain{sur} $L^\infty(]0, \omega[)$ \uncertain{se} \uncertain{correspondent} \uncertain{par} $W(\exp f)$ \ill{} $P$ \uncertain{une} \uncertain{fonction} \uncertain{de} \uncertain{Weyl} \uncertain{sur} $L^\infty(]0, \omega[)$ \ill{} $P(\exp f)$ \ill{}, \uncertain{soit} $P(f) = $ \ill{} \uncertain{si} $f$ \uncertain{est} \uncertain{une} \uncertain{fonction} \ill{} $P(f)$ : \struck{\ill{}} \ill{} $\varphi_f$ \ill{} $0 \leq s \leq \omega$ \ill{} \uncertain{de} $W(\exp f)$, i.e. \uncertain{de} \ill{} $\varphi_f$, \uncertain{si} \ill{}, \uncertain{de} $f$ \ill{} \struck{\ill{}} \uncertain{fonction} : \ill{} \uncertain{croissante}. \uncertain{Alors} \ill{}
\[
  W(\exp f) = \sup_i \Bigl( a_i + \int_0^{\omega} \varphi_i\,\varphi_f \Bigr) \qquad
  \uncertain{P}(f) = \sup_i \Bigl( a_i + \int_0^{\infty} \varphi_i \log\varphi_f \Bigr)
\]
\note{le premier membre de la première formule est écrit « $P(\exp f)$ » avec un $P$ que la ligne suivante fait lire $W$ ; le $\varphi_f$ de la première porte un exposant surchargé} \ill{} \uncertain{fonction} \struck{\uncertain{positive}} \uncertain{telle} \uncertain{que} \ill{} $\varphi_i(s) > 0$. \uncertain{Le} \ill{} \uncertain{définition} \uncertain{des} $W$ \ill{} \uncertain{continues} \ill{} \uncertain{semi-continuité} \uncertain{pour} \ill{}, \uncertain{représentation} \uncertain{...} \uncertain{...} \uncertain{fonction} \uncertain{de} \uncertain{Pólya} \uncertain{sur} \uncertain{les}

\page{13}
\uncertain{parties} \uncertain{positives} \uncertain{de} $L^\infty(]0, \omega[)$ \ill{} \uncertain{fonction} \uncertain{de} Weyl. \note{en interligne, cerné : « (\uncertain{ou} \uncertain{des} \ill{} \uncertain{ci-dessous}) »}

\emph{Exemples de fonctions de Weyl}. Soit $\alpha$ une \uncertain{fonction} \uncertain{sur} $\mathbf{R}^{+}$, \struck{convexe}, \uncertain{nulle} \uncertain{à} \uncertain{l'origine} \ill{} \struck{Alors $\int \alpha(\varphi)$} \uncertain{telle} \uncertain{que} $\alpha(\exp s)$ soit \uncertain{fonction} \uncertain{convexe} \uncertain{de} $s$ (i.e. $\alpha(s^\lambda t^{(1-\lambda)}) \leq \alpha(s)^\lambda\alpha(t)^{1-\lambda}$). Alors
\[
  W_\alpha(f) = \int_0^\omega \alpha(\varphi_f) \struck{\ill{}}
\]
\uncertain{est} \uncertain{une} \uncertain{fonction} \uncertain{de} Weyl, \uncertain{croissante} \uncertain{si} \uncertain{et} \uncertain{seulement} \uncertain{si} $\alpha$ \uncertain{est} \uncertain{croissante}. \uncertain{En} \uncertain{particulier} \uncertain{prenant} $\alpha_p(s) = s^p$ ($0 < p < +\infty$) \uncertain{on} \uncertain{trouve}
\[
  \text{(1)}\qquad W_{\alpha_p}(f) = N_p(f)^p = \int |f|^p \struck{\ill{}}
\]

\emph{Extension d'une fonction de Weyl croissante} \uncertain{Une} \uncertain{fonction} \uncertain{de} Weyl \uncertain{croissante} \uncertain{sur} $(L^\infty)^{+}(]0, \omega[)$ \uncertain{est} \uncertain{considérée} \uncertain{comme} \uncertain{restriction} \uncertain{d'une} \uncertain{fonction} \uncertain{de} Weyl \uncertain{croissante} \uncertain{sur} $(K^\infty)^{+}(]0, +\infty[)$. \struck{Soit $W$ une fonction de Weyl croissante \uncertain{sur} \ill{}} Soit $W$ \uncertain{une} \uncertain{telle} \uncertain{fonction} \uncertain{de} Weyl. Soit $f$ \uncertain{une} \uncertain{fonction} \uncertain{mesurable} \struck{\ill{}} positive sur un \uncertain{espace} \uncertain{mesuré} $M$, \uncertain{on} \struck{pose} \uncertain{pose} \uncertain{alors}
\[
  W(f) = \sup W(\varphi_g)
\]
\note{sous le $W$ du second membre, un indice raturé ; le $\varphi_g$ est surchargé} \struck{\ill{}} \uncertain{la} \uncertain{borne} \uncertain{supérieure} \uncertain{étant} \uncertain{prise} \uncertain{sur} \uncertain{les} \uncertain{fonctions} $g$ \ill{} \uncertain{telles} \uncertain{que} $0 \leq g \leq f$. \uncertain{On} \uncertain{a} \struck{\ill{}} $W(f) = W(\varphi_f)$ \note{cerné d'un trait} \ill{} \uncertain{si} $|f|$ \ill{} \uncertain{et} \uncertain{si} $|f|$ \ill{} \struck{\uncertain{définition}} $W(f) = W(\varphi_{|f|})$ \note{cerné d'un trait} \uncertain{L'ensemble} \uncertain{des} \uncertain{fonctions} $f$ \uncertain{telles} \uncertain{que} $W(f) < +\infty$, \ill{} \uncertain{dérivé}, \uncertain{et} $W$ \uncertain{y} \uncertain{est} \ill{}

\page{14}
\uncertain{topologie} (\uncertain{faible}) \uncertain{dans} \uncertain{les} \uncertain{parties} \uncertain{compactes} \uncertain{de} $L^1(K)$, \ill{} \struck{\ill{} \uncertain{parties} \uncertain{compactes} \uncertain{de} $L^\infty(K)$}. \uncertain{i.e.} $\sigma(L^\infty(K), L^\infty(K))$ \note{sic, le second $L^\infty$ est surchargé}, \uncertain{donc} (\uncertain{dans} \uncertain{les} \uncertain{parties} \uncertain{bornées} \uncertain{de} $L^\infty(K)$), \uncertain{pour} $\sigma(L^\infty(K), L^1(K))$. \note{une accolade en marge gauche embrasse ces deux lignes}

\marginal{en diagonale dans la marge gauche, sur six lignes : \ill{} \uncertain{voisinages} \ill{} \uncertain{voisinage} \ill{}}

\uncertain{Considérons} \uncertain{maintenant}, \uncertain{dans} $\underline{\mathbf{R}} \times L^\infty(K)$, \uncertain{l'ensemble} $A_P(K)$ \note{l'indice de $A$ est surchargé} \uncertain{des} \uncertain{couples} $(\lambda, f)$ tels que $\lambda \geq P(f)$. $P_2$ \uncertain{signifie} \uncertain{que} $A_P$ \uncertain{est} \uncertain{convexe}, $P'_3$ \uncertain{que} $A_P$ \uncertain{est} \uncertain{fermé}, \uncertain{donc} (\uncertain{Hahn}-\uncertain{Banach}) $A_P$ \uncertain{est} \uncertain{intersection} \uncertain{de} \uncertain{demi-espaces} \uncertain{fermés}, \uncertain{i.e.} \uncertain{de} \uncertain{demi-espaces} \ill{} $\lambda \geq L_i(f) + a_i$ \note{le $\lambda$ est surchargé} ($L_i$ \uncertain{forme} \uncertain{linéaire} \uncertain{continue} \uncertain{sur} $L^\infty(K)$, i.e. $L_i \in L^1(K)$, \struck{\ill{}} $a_i \in \underline{\mathbf{R}}$) \uncertain{d'où} \struck{$P(f) = \ill{}$}
\[
  P(f) = \sup_i \bigl( L_i(f) + a_i \bigr).
\]
En vertu de $P_1$ on a \struck{\ill{}} aussi, \uncertain{pour} \uncertain{tout} $i$ \uncertain{et} \uncertain{tout} $g \in L^\infty(K)$ \note{ajouté au-dessus} \uncertain{spectralement} \uncertain{équivalent} \uncertain{à} $f$,
\[
  P(f) \geq L_i(g) + a_i
\]
\uncertain{d'où} \uncertain{en} \uncertain{vertu} \uncertain{de} \uncertain{prop.}~2, \uncertain{prise} \uncertain{sur} \uncertain{les} $g$ \note{deux lignes cernées d'un trait} :
\[
  P(f) \geq \int \varphi_{L_i}\varphi_f + a_i, \qquad \text{Posons } \varphi_{L_i} = \varphi_i :
\]
\[
  \text{(1)}\qquad \boxed{\ P(f) = \sup_{i \in I}\ a_i + \int_0^{\omega} \varphi_i\,\varphi_f\ }
\]
\note{la borne supérieure $\omega$ de l'intégrale est surchargée ; à droite, un mot raturé}

\page{15}
\note{Feuillet plus petit ; une longue diagonale barre sa moitié inférieure, de (2) jusqu'au bas, sans que l'on puisse dire si elle annule le passage}

\uncertain{les} $\varphi_i$ \uncertain{sont} \uncertain{des} \uncertain{fonctions} \uncertain{décroissantes} \note{ajouté en interligne, cerné : « \ill{} \uncertain{sur} $]0, |K|]$ »} (\uncertain{par} \ill{} \uncertain{bornée} \ill{} \uncertain{de} $K$). \uncertain{Réciproquement}, \uncertain{si} $P$ \uncertain{est} \uncertain{définie} \uncertain{sur} $L^\infty(K)$ \uncertain{par} \uncertain{une} \uncertain{formule} (1), $P$ \uncertain{est} \uncertain{une} \uncertain{fonction} \uncertain{de} \uncertain{Pólya}. \ill{} $P$ \uncertain{satisfaisant} \ill{} $P$ \uncertain{est} \emph{croissante}.

\marginal{en marge gauche, sur deux colonnes sideways, de la même encre. Première colonne : \uncertain{Soit} \ill{} \uncertain{de} \uncertain{Pólya}. \uncertain{Alors} $P$ \ill{} ($0 < \omega < +\infty$) \ill{} $P$ \ill{} $|K|$ \ill{} \uncertain{fonctions} \ill{} \emph{Proposition 3} \uncertain{Pour} \uncertain{que} $P$ \uncertain{sur} $]0, \omega[$ \ill{} \uncertain{fonction} \uncertain{de} \uncertain{Pólya}, \ill{} (1), \ill{} (\uncertain{fonctions} \uncertain{décroissantes} $\in L^1(]0, \omega[)$, \ill{}) \ill{} \uncertain{croissante} \ill{} $\varphi_i$ \ill{}. Seconde colonne : \emph{Définition} \uncertain{Une} \uncertain{fonction} \uncertain{de} \uncertain{Pólya} \ill{} \uncertain{est} \uncertain{dite} \ill{} : \ill{} \struck{\ill{}} \ill{}}

\uncertain{Alors} \uncertain{on} \uncertain{voit} \uncertain{aussitôt} \ill{} \uncertain{fonctions} (\uncertain{et} \uncertain{réciproquement}) \ill{} \uncertain{on} \uncertain{peut} \uncertain{prendre} \uncertain{les} $\varphi_i$ \uncertain{dans} (1), \ill{} \note{cerné d'un trait, avec une ligne biffée} \uncertain{nulles} \uncertain{sur} $[|K|, +\infty[$, \uncertain{et} \struck{\ill{}} \ill{} (\ill{} \uncertain{de} $K$) \ill{} \uncertain{définie} \ill{}
\[
  \text{(2)}\qquad \boxed{\ P(f) = \sup_i\ a_i + \int_0^{\infty} \varphi_i\,\varphi_f\ }
\]
\uncertain{où} \uncertain{les} $\varphi_i$ \uncertain{sont} \uncertain{des} \uncertain{fonctions} \uncertain{décroissantes} \uncertain{sur} $\mathbf{R}^{+}$, \struck{\ill{}} \note{une ligne biffée, cernée avec les mots « \uncertain{sommables} » au-dessus} \ill{} \uncertain{et} \ill{}, \ill{} (2) \uncertain{peut} \uncertain{s'écrire}
\[
  \text{(3)}\qquad P(f) = P^{*}(\varphi_f)
\]
\uncertain{où} $P^{*}$ \uncertain{est} \uncertain{la} \uncertain{fonction} \uncertain{de} \uncertain{Pólya} \uncertain{sur} $\mathbf{R}^{+}_{*}$ \uncertain{donnée} \uncertain{par}
\[
  \text{(4)}\qquad P^{*}(\varphi) = \sup_i \Bigl( a_i + \int_0^{\infty} \varphi_i\,\varphi \Bigr)
\]

\page{16}
Si \uncertain{par} \uncertain{exemple} \ill{} \uncertain{pour} \uncertain{une} \ill{} $\omega$ \ill{} (*), \struck{$W(f) < +\infty$ \ill{} $f \in L^p(M)$, \ill{}} \uncertain{on} \uncertain{a} $W(f) = \int_M |f|^p$, \uncertain{donc} $W(f) < +\infty$ \uncertain{équivaut} \uncertain{à} $N_p(f) < +\infty$.

\section{Inégalités liées aux fonctions de Weyl}
\note{Titre de sa main, page 16 : « 3. Inégalités \uncertain{liées} \uncertain{aux} \struck{fonctions de Weyl} fonctions de Weyl » — les mots biffés sont réécrits au-dessus. C'est le « N° 3 » du plan de la page 2.}

Soit $0 < \omega < +\infty$.

\emph{Prop 4} Soient $f, g$ \uncertain{deux} \uncertain{fonctions} \uncertain{décroissantes} \uncertain{sur} $]0, \omega[$ (\uncertain{localement} \uncertain{sommables} \ill{}), \uncertain{telles} \uncertain{que}
\[
  \text{(1)}\qquad \boxed{\ \int_0^t f \leq \int_0^t g\ } \qquad\text{pour } 0 \leq t \leq \omega
\]
\uncertain{Alors} \uncertain{pour} \uncertain{toute} \uncertain{fonction} \uncertain{de} Pólya \add{\uncertain{croissante}} \note{écrit au-dessus} $P$ \uncertain{sur} $L^\infty(]0, \omega[)$, \uncertain{on} \uncertain{a}
\[
  \text{(1')}\qquad P(f) \leq P(g)
\]
\note{en marge, un « ① » cerclé, avec une accolade sur les trois lignes qui suivent} (\uncertain{On} \uncertain{obtient} \uncertain{une} \uncertain{réciproque} \ill{} \uncertain{d'ailleurs} \uncertain{trivialement}, \uncertain{à} \uncertain{savoir} \uncertain{que} \uncertain{pour} \uncertain{tout} $t$, $f \to \int_0^t \varphi_f$ \uncertain{est} \uncertain{une} \uncertain{fonction} \uncertain{de} \uncertain{Pólya} \uncertain{croissante} \uncertain{sur} $L^\infty(]0, \omega[)$).

En vertu de N°~2, (1), il \uncertain{suffit} \uncertain{de} \uncertain{prouver} \uncertain{que}
\[
  \int_0^\omega f\varphi \leq \int_0^\omega g\varphi
\]
\uncertain{pour} \uncertain{toute} \uncertain{fonction} \emph{positive} \emph{décroissante} \uncertain{sommable} \uncertain{sur} $(0, \omega)$. \uncertain{Or} \uncertain{cela} \uncertain{résulte} \uncertain{facilement} \uncertain{de} (1) \uncertain{par} \uncertain{intégration} \uncertain{par} \uncertain{parties}, \uncertain{en} \uncertain{posant} $F(t) = \int_0^t f$, \ill{} :
\[
  \text{(3)}\qquad \int_0^\omega f\varphi = [\varphi F]_0^\omega + \int_0^\omega F(-d\varphi) = \varphi(\omega)F(\omega) + \int_0^\omega F(-d\varphi)
\]
\note{le numéro (2) n'apparaît pas ; une longue accolade en marge gauche embrasse la page de (1') jusqu'ici}

\emph{Remarque} \note{suivi d'un « ① » cerclé, relié d'une flèche à la ligne} \uncertain{On} \uncertain{voit} \uncertain{sur} (3) \uncertain{que} \uncertain{si} $F(\omega) = G(\omega)$, i.e. $\int_0^\omega f = \int_0^\omega g$, \uncertain{il} \uncertain{n'est} \uncertain{pas} \uncertain{nécessaire} \uncertain{de} \uncertain{supposer} $\varphi$ \uncertain{positive} \uncertain{dans} \uncertain{la} \uncertain{conclusion}, \uncertain{de} \uncertain{sorte} \uncertain{que} \uncertain{la} \uncertain{Prop.}~4 \ill{} \uncertain{pour} \uncertain{toute} \uncertain{fonction} \uncertain{de} \uncertain{Pólya} \ill{} \uncertain{non} \uncertain{nécessairement} \uncertain{croissante} \ill{} \uncertain{des} \ill{}

\page{17}
\ill{} \uncertain{inégalités} \ill{} \uncertain{dans} \uncertain{les} \ill{} \uncertain{avec} \ill{} \uncertain{un} \ill{} (\ill{}).

\emph{Corollaire 1} Soient $f, g$ \struck{\ill{} $(K^\infty)^{+}(]0, \omega[)$ \ill{}} \uncertain{fonctions} \uncertain{positives} \ill{} \uncertain{sur} $M$ \note{« sur $M$ » écrit au-dessus, avec un « $N$ » raturé} \uncertain{telles} \uncertain{que}
\[
  \text{(4)}\qquad \Delta_f(t) \leq \Delta_g(t) < +\infty \qquad 0 \leq t \leq \omega
\]
\uncertain{Alors} \uncertain{pour} \uncertain{toute} \uncertain{fonction} \uncertain{de} Weyl \struck{\uncertain{croissante} \uncertain{sur} \ill{}} \uncertain{croissante} \uncertain{sur} $(K^\infty)^{+}(]0, \omega[)$ \uncertain{on} \uncertain{a}
\[
  \text{(5)}\qquad W(f) \leq W(g)\,.
\]
\struck{\ill{} \uncertain{sur} \ill{}} \uncertain{Quitte} \uncertain{à} \uncertain{remplacer} \ill{} \uncertain{on} \uncertain{peut} \uncertain{supposer} $M = ]0, \omega[$, \struck{$f$ et $g$ \ill{}} \ill{} \note{une ligne cernée d'un trait : « \uncertain{puis} \ill{} »} \struck{$f$ et $g$ \uncertain{décroissantes}} \ill{}. \uncertain{Quitte} \uncertain{à} \ill{} \ill{} $0 < \omega < +\infty$, \uncertain{puis} \struck{\ill{}} \ill{} $f$ \ill{} $\exp f$, \uncertain{et} \ill{} \uncertain{fonctions} \uncertain{de} \uncertain{Weyl} \uncertain{croissantes} \ill{} \note{une ligne biffée} \ill{} \uncertain{fonctions} \uncertain{de} \uncertain{Pólya} \uncertain{croissantes} \ill{} \uncertain{N°}~2 \ill{} \uncertain{Prop.}~4. \struck{\ill{}} CQFD
\note{le reste de la page est blanc ; « CQFD » est souligné}

\section{Mesures spectrales et réarrangements spectraux d'opérateurs}
\note{Titre de sa main, page 18 : « 4 Mesures spectrales et réarr\add{angements} \uncertain{spectraux} d'opérateurs. » C'est le « N° 4 » du plan de la page 2. Dans ces trois pages, une main postérieure à l'encre noire corrige $\Omega$ en $M$ à chaque occurrence ; on transcrit $M$ et on signale la première correction.}

\page{18}
\uncertain{Dans} \uncertain{tout} \uncertain{ce} \uncertain{numéro}, $\mathcal{A}$ \uncertain{est} \uncertain{une} \uncertain{algèbre} \uncertain{de} \uncertain{von} Neumann (\uncertain{algèbre} \uncertain{involutive} \uncertain{faiblement} \uncertain{fermée} d'opérateurs \uncertain{dans} \uncertain{un} \uncertain{espace} \uncertain{hilbertien} \ill{}). \struck{\uncertain{Trace} \uncertain{normale} (} \uncertain{Une} \uncertain{trace} \uncertain{sur} $\mathcal{A}$ \uncertain{est} \uncertain{une} \struck{\ill{}} \uncertain{fonction} $T$ \uncertain{sur} \uncertain{le} \uncertain{cône} \uncertain{positif} $\mathcal{A}^{+}$ de $\mathcal{A}$, à valeurs $0 \leq T(A) \leq +\infty$, \struck{additive, positivement homogène} \uncertain{telle} \uncertain{que}
\[
  T(A+B) = T(A) + T(B) \quad\text{et}\quad T(\lambda A) = \lambda T(A)
\]
($A, B \in \mathcal{A}^{+}$, $\lambda \geq 0$ \uncertain{scalaire} ; \uncertain{on} \uncertain{pose} $0\cdot(+\infty) = 0$), et \uncertain{invariante} \uncertain{par} \uncertain{les} \uncertain{transformations} \uncertain{unitaires}. \uncertain{Alors} l'ens. des $A \in \mathcal{A}^{+}$ tels que $T(A) < +\infty$ \uncertain{est} \uncertain{le} \uncertain{cône} \uncertain{positif} $\underline{a}^{+}$ \note{ajouté au-dessus} \uncertain{d'un} \uncertain{idéal} \uncertain{bilatère} \uncertain{involutif} $\underline{a}$, \uncertain{et} $T$ \uncertain{se} \uncertain{prolonge} \uncertain{en} \uncertain{une} \uncertain{forme} \uncertain{linéaire} \uncertain{sur} $\underline{a}$, \uncertain{notée} \uncertain{encore} $T$. On a
\[
  T(AB) = T(BA) \qquad A \in \underline{a},\ B \in \mathcal{A}
\]
$T$ \uncertain{est} \uncertain{dite} \emph{\uncertain{semi-finie}} \uncertain{si} \ill{} \uncertain{pour} \uncertain{tout} \struck{$A \in \mathcal{A}^{+}$, $A \neq 0$, \uncertain{existe}} $A \in \mathcal{A}$, $A > 0$, \uncertain{existe} $B$ \uncertain{tel} \uncertain{que} $0 < B < A$ \uncertain{et} $T(B) \struck{\neq} < +\infty$. $T$ \uncertain{est} \uncertain{dite} \uncertain{normale} \uncertain{si}
\[
  T(\textstyle\sup_i A_i) = \sup_i T(A_i)
\]
\uncertain{pour} \uncertain{toute} \uncertain{famille} \uncertain{filtrante} \uncertain{croissante} \uncertain{majorée} $(A_i)$ \uncertain{dans} $\mathcal{A}^{+}$. \uncertain{Si} $T$ \uncertain{est} \uncertain{normale}, \uncertain{la} \uncertain{définition} \ill{} \uncertain{et} \uncertain{seulement} \uncertain{si} \ill{}. \uncertain{Enfin} $T$ \uncertain{est} \uncertain{dite} \emph{\uncertain{fidèle}} \uncertain{si} \struck{$T(A)$ $A \in \mathcal{A}^{+}$} $A \geq 0$, $T(A) = 0 \Rightarrow A = 0$ \ill{} $\mathcal{A}$ \ill{}. \uncertain{Dans} \uncertain{la} \uncertain{suite} \ill{} \uncertain{trace} \uncertain{normale} \uncertain{semi-finie} \uncertain{fidèle} $T$ \note{cerné d'un trait}, \uncertain{dont} \uncertain{l'idéal} \uncertain{de} \uncertain{définition} \uncertain{est} $\underline{a}$ \ill{} \uncertain{idéal} \uncertain{bilatère} \ill{} \uncertain{définition} \uncertain{de} $\underline{a}$. Soit $\underline{b}$ l'adhérence \uncertain{normique} \ill{} \uncertain{de} $\underline{a}$, \uncertain{c'est} \uncertain{un} \uncertain{idéal} \uncertain{bilatère} \uncertain{involutif} \uncertain{fermé}

\page{19}
\note{Deux longues diagonales barrent la page entière, du haut à droite au bas à gauche, et plusieurs blocs y sont en outre cernés ou biffés ligne à ligne ; le passage est annulé et repris à la page 20. On transcrit ce qui se lit.}

\struck{\ill{}} \uncertain{tel} \uncertain{que} $\underline{b}$ \uncertain{soit} \uncertain{un} \ill{}, i.e. \uncertain{pour} \ill{} \struck{\ill{} (i.e. \ill{} \uncertain{dans} $\mathcal{A}$)} \uncertain{pour} \uncertain{que} $\underline{b} = \mathcal{A}$, il \uncertain{faut} \uncertain{et} \uncertain{il} \uncertain{suffit} \uncertain{que} $\underline{a} = \mathcal{A}$, \uncertain{et} i.e. que $T$ \uncertain{soit} \uncertain{une} \uncertain{trace} \emph{\uncertain{finie}} (i.e. $T(A) < \infty$ \uncertain{pour} \uncertain{tout} $A \in \mathcal{A}^{+}$).

\struck{\emph{Proposition 1} Soit $A \in \mathcal{A}$ \ill{}} \uncertain{Alors} \uncertain{les} \uncertain{propriétés} \uncertain{spectrales} \ill{} \uncertain{de} $\underline{b}$ \ill{} \struck{En \uncertain{fait}, \uncertain{les} \uncertain{propriétés} \uncertain{spectrales} \uncertain{de} $\underline{b}$ \ill{} \uncertain{déduisent}} \uncertain{Soit} $E$ \uncertain{un} \uncertain{projecteur} \ill{} \uncertain{dans} $\underline{a}$ \ill{} \uncertain{pour} \uncertain{qu'un} \uncertain{projecteur} \uncertain{de} $\mathcal{A}$ \uncertain{soit} \uncertain{dans} $\underline{a}$, il \uncertain{faut} \uncertain{et} \uncertain{il} \uncertain{suffit} \uncertain{que} $T(E) < +\infty$. Or \ill{}

\struck{\emph{Proposition 1} Soit $A$ \uncertain{un} \uncertain{opérateur} \uncertain{normal} \uncertain{de} $\underline{b}$, \ill{} \uncertain{hermitien} \ill{} $\underline{C}$ \ill{} \uncertain{de} \ill{} $\underline{a}$. \uncertain{Alors} $f(A) \in \underline{a}$ \uncertain{pour} \uncertain{toute} \uncertain{fonction} \ill{} \struck{\ill{}} \ill{} $\underline{a} \subset \underline{b}$, $\subset$ \ill{} \uncertain{sur} $A$ : \uncertain{un} \ill{} $C_0(\Omega)$. \ill{} \uncertain{définit} \uncertain{une} \ill{} \uncertain{fonction} \ill{} \uncertain{sur} $C_0(\Omega)$. \uncertain{Ceci} \uncertain{...} \uncertain{partie} \uncertain{positive} \ill{} \uncertain{séparations} \ill{} \uncertain{caractéristique} \ill{} \uncertain{important}. \uncertain{Il} \ill{}} \note{ce bloc, cerné et barré, contient aussi : « \struck{\emph{Corollaire 1}} $A$ \ill{} $C_0(\Omega)$, $A$ \ill{} $f(A)$ \uncertain{est} \uncertain{l'opérateur} \ill{} $f(\varphi)$ \uncertain{sur} \ill{} \uncertain{borné} \ill{} \uncertain{bornée} \ill{} \uncertain{ensemble} \ill{} $E$ \uncertain{est} \uncertain{l'application} \ill{} \uncertain{et} \uncertain{soit} \uncertain{donc} \ill{} \uncertain{pour} \ill{} »}

\emph{Corollaire} \struck{Soit} $K(\Omega) \subset \underline{a}$
\note{le reste de la page est blanc}

\page{20}
\emph{Proposition 1} \note{le numéro est surchargé} Soit $\underline{a}$ \uncertain{un} \uncertain{idéal} \uncertain{bilatère} \ill{} \uncertain{de} \uncertain{l'algèbre} \uncertain{de} \uncertain{von} Neumann $\mathcal{A}$, $\underline{b}$ \uncertain{son} \uncertain{adhérence} \uncertain{normique} \ill{} $E \in \underline{b}$ \struck{\uncertain{fonction} \uncertain{hermitienne} \ill{} $\underline{C}$} \ill{} \uncertain{sous-algèbre} \ill{} \uncertain{de} $\underline{b}$, \uncertain{identifiée} \uncertain{à} \uncertain{une} \struck{\uncertain{fonction} \ill{}} \ill{} $K(\struck{\Omega}\add{M}) \subset \underline{a}$ \note{le $\Omega$ biffé et corrigé en $M$ à l'encre noire, ici et dans la suite}. \uncertain{Si} $\underline{C}$ \uncertain{est} \ill{}, \uncertain{alors} \ill{} \uncertain{fonction} \ill{} $K$ \uncertain{de} $\Omega$ \ill{} \uncertain{fonction} \uncertain{hermitienne} \ill{} $\underline{C}$ \uncertain{est} \uncertain{une} \uncertain{sous-algèbre} \ill{} \uncertain{Alors} \ill{} \uncertain{pour} \ill{} $E \in \underline{b}$ \ill{} $\underline{a}$. \struck{\ill{}} \uncertain{soit} $\underline{c}$ \ill{} $\underline{C}$, \uncertain{borné} \ill{} \uncertain{opérateurs} \ill{} \uncertain{relation} \uncertain{compact} \ill{} \ill{} $C_0(\struck{\Omega}\add{M})$, \uncertain{alors} \struck{\ill{}} \ill{} $\underline{C}$ \uncertain{est} \uncertain{unitairement} \ill{}, \uncertain{alors} \ill{} \uncertain{contenu} \uncertain{dans}, \uncertain{un} \ill{} \uncertain{compact} \ill{} $\Omega$ \uncertain{est} \struck{\ill{}} \ill{} \uncertain{d'après} \ill{} \struck{\ill{}} $f \in C_0(\struck{\Omega}\add{M})$ \ill{} \struck{\ill{}} \note{trois lignes biffées, dont une à l'encre noire}

\marginal{en marge gauche, sideways, à hauteur de la démonstration : \ill{} \uncertain{fermé} \ill{} \uncertain{idéal} \ill{} $\underline{b}$ \ill{} $\underline{a}$}

il \uncertain{existe} \uncertain{un} \uncertain{projecteur} \ill{} \struck{\ill{}} \uncertain{qui} \struck{\ill{}} \uncertain{Soit} $E = \lim A_i$ ($A_i \in \underline{a}$), \uncertain{d'où} $E = \lim A_i^{*}A_i$, \uncertain{donc} $E = \lim B_i$ ($B_i \geq 0$, $B_i \in \underline{a}$) \uncertain{d'où} $E = \lim E B_i E$, i.e. $E = \lim C_i$, \uncertain{où} $C_i \in \underline{a}$, $E C_i E = C_i$. \uncertain{On} \uncertain{a} \ill{} $C_i \leq \|C_i\| E$, \uncertain{d'où} \note{en interligne : « $\dfrac{C_i}{\|C_i\|} \leq E$ »} \ill{} \ill{} \emph{$E = \lim_i D_i$}, \uncertain{avec} $D_i \in \underline{a}$, $0 \leq D_i \leq E$. \uncertain{Mais} \uncertain{si} $\|E - D_i\| \leq \frac{1}{2}$, \uncertain{alors} $E \leq 2 D_i$ \note{le $D$ est surchargé} (\uncertain{d'après} \ill{}) \uncertain{d'où} $E \in \underline{a}$.

\uncertain{Pour} 2 \uncertain{démontrer} \ill{} \struck{\ill{} \uncertain{maximale} \ill{} \uncertain{de} $\mathcal{A}$)} \uncertain{si} $f \in K^{+}(\struck{\Omega}\add{M})$ \ill{} \uncertain{caractéristique} \uncertain{d'un} \ill{} \uncertain{et} \ill{} \uncertain{note} \uncertain{que} \ill{} \struck{\ill{}} \ill{} $0 \leq f \leq \varphi_K \leq g$, \uncertain{fonction} \ill{} \note{au-dessus : « $\varphi_K$ »} \ill{} $K \subset \struck{\Omega}\add{M}$, $g \in K^{+}(\struck{\Omega}\add{M})$. \uncertain{Comme} $U_{\varphi_K} \leq U_g \in \underline{b}$, \uncertain{d'où} $U_{\varphi_K} \in \underline{b}$, \uncertain{d'où} \ill{} \uncertain{ce} \uncertain{qui} \uncertain{précède} $U_{\varphi_K} \in \underline{a}$, \uncertain{d'où} $U_f \in \underline{a}$. \uncertain{Comme} \ill{} \ill{}, (\uncertain{c'est} \uncertain{d'ailleurs} \ill{} \uncertain{fonction} \ill{} \uncertain{de} \uncertain{l'algèbre} \uncertain{de} $\mathcal{A}$), \struck{\uncertain{soit} $K$ \ill{} \uncertain{compact} \ill{} $U_{\varphi_K} \in \underline{a}$ \ill{}}
\note{le lot s'arrête ici ; la démonstration se poursuit au feuillet suivant}

\end{document}
